% Le resume presente rapidement le probleme etudie et l'approche utilisee.
% Une demi-page au grand maximum.

Une piscine couverte s'évacue en urgence : les baigneurs sortent de l'eau par les échelles et
le rebord, traversent une plage dont le carrelage est mouillé, puis se pressent devant une
porte unique. Nous cherchons comment cet enchaînement altère le diagramme fondamental
débit--densité de la foule, et lequel de deux mécanismes l'emporte : un effet tampon, où la
traversée désynchronise les arrivées à la porte, ou un effet inertiel, où le sol glissant
empêche les baigneurs de réguler leur approche.

Nous utilisons un modèle particulaire à exclusion dure : à chaque pas, un déplacement libre
est projeté sur l'ensemble des positions admissibles, et la vitesse est relue sur le
déplacement réalisé. Les trois milieux (eau, carrelage mouillé, sol sec) diffèrent par la
vitesse souhaitée, le temps de relaxation et l'accélération maximale que permet l'adhérence.
Un guidage par champ eikonal mène chacun vers la sortie, et une régulation vitesse--espacement,
étalonnée sur les mesures de Seyfried et al.\ au goulot, fixe le débit.

Dans la piscine de référence, le carrelage mouillé ne change pas l'évacuation de façon
mesurable. Quand un sol très glissant touche la file d'une porte saturée, l'effet inertiel
apparaît, sur vingt graines : la porte perd jusqu'à \SI{19}{\percent} de son débit, et la file
devient plus dense et plus lente à densité égale. L'effet tampon, lui, n'a pas été mis en
évidence.
